% =============================================================================
% 1.11 Solutions existantes et synthèse de l'état de l'art
% =============================================================================
\section{Solutions existantes et synthèse}
\label{sec:ch1-solutions}

\subsection{Solutions du marché}

Le marché de la gestion des appels d'offres et de la gestion de projet BTP
propose plusieurs familles d'outils. Le tableau~\ref{tab:solutions-marche}
en donne un aperçu représentatif, sans prétendre à l'exhaustivité.

\begin{btptable}{Familles de solutions existantes}
\label{tab:solutions-marche}
\begin{tabularx}{\linewidth}{@{}L{2.8cm} Y C{1.6cm} C{2.0cm}@{}}
\toprule
\rowcolor{btppale}
\btpth{Famille} & \btpth{Fonctionnalités typiques} & \btpth{IA générative} & \btpth{Capitalisation} \\
\midrule
Gestion d'appels d'offres & Suivi des consultations, collaboration, modèles de réponse. & Limitée & Passive \\
Gestion de projet BTP & Planning, coûts, documents de chantier. & Absente & Passive \\
Veille des marchés publics & Détection d'opportunités, alertes. & Absente & Sans objet \\
GED généraliste & Archivage, circuits, recherche plein texte. & Émergente & Passive \\
\bottomrule
\end{tabularx}
\end{btptable}

Ces solutions partagent quatre limites, déjà identifiées à la
section~\ref{sec:ch1-ged} et que l'état de l'art technologique permet
désormais d'adresser~: recherche sémantique absente ou rudimentaire,
absence de modélisation des relations, capitalisation réduite à
l'archivage, et génération automatique inexistante ou non sourcée.

Aucune, à notre connaissance, n'articule dans un même système une ingestion
structurelle des pièces d'un DAO, une recherche hybride sur un patrimoine
d'entreprise, une génération section par section avec citations, et une
supervision humaine obligatoire avant export --- le tout dans le cadre
juridique des marchés publics camerounais.

\subsection{Tendances émergentes}

Quatre tendances se dessinent néanmoins dans la littérature et chez les
éditeurs~:

\begin{enumerate}
  \item l'usage des grands modèles de langage pour produire un premier
        brouillon de réponse~;
  \item l'extraction automatique des exigences d'un dossier entrant~;
  \item le RAG comme architecture d'aide à la rédaction, plutôt que comme
        chatbot générique~;
  \item l'orchestration de workflows couvrant le processus de bout en bout.
\end{enumerate}

Ces tendances confirment la pertinence de notre problématique. Elles en
confirment aussi le risque~: une génération non contrainte, sans citation
ni revue, n'est pas une aide, c'est un nouveau vecteur d'irrégularité.

\subsection{Synthèse de l'état de l'art}

Le tableau~\ref{tab:synthese-eoa} croise la maturité des domaines parcourus
et leur pertinence pour la réponse aux appels d'offres.

\begin{btptable}{Maturité des domaines et pertinence pour les appels d'offres}
\label{tab:synthese-eoa}
\begin{tabularx}{\linewidth}{@{}L{4.4cm} C{3.0cm} Y@{}}
\toprule
\rowcolor{btppale}
\btpth{Domaine} & \btpth{Maturité} & \btpth{Pertinence} \\
\midrule
Traitement automatique du langage & \btpMature & Très élevée \\
Grands modèles de langage & \btpMature{} (évolutive) & Très élevée \\
Recherche d'information lexicale (BM25) & \btpMature & Très élevée \\
Fusion de classements (RRF) & \btpMature & Très élevée \\
Génération augmentée par récupération & \btpMature & Très élevée \\
Bases vectorielles & \btpMature & Très élevée \\
Compréhension de documents & \btpEnDev & Très élevée \\
Résolution d'entités & \btpMature & Élevée \\
Graphes de connaissances & \btpEnDev & Élevée en seconde phase \\
Agents fondés sur les LLM & \btpEmergent & Ponctuelle, non centrale \\
Supervision humaine dans la boucle & \btpMature & Constitutive \\
\bottomrule
\end{tabularx}
\end{btptable}

De ce panorama se dégagent cinq lacunes que le présent travail entend
combler~:

\begin{enumerate}
  \item l'absence de solution intégrée couvrant le cycle analyse--rédaction
        --validation--export d'un dossier d'offre~;
  \item le manque de traçabilité des sources dans les approches génératives
        existantes~;
  \item une capitalisation encore passive, réduite à l'archivage~;
  \item l'absence de supervision humaine comme étape obligatoire, et non
        comme option~;
  \item la non-adaptation au cadre juridique et documentaire des marchés
        publics camerounais.
\end{enumerate}

\begin{btpbox}[Ce que le chapitre suivant doit trancher]
L'état de l'art justifie une architecture de génération augmentée par
récupération, adossée à une recherche hybride, à une ingestion
structurelle et à une revue humaine. Il ne dit pas \emph{comment}
articuler ces briques, ni selon quels critères les évaluer. C'est l'objet
du chapitre~\ref{chap:methodologie}.
\end{btpbox}

\btpEndMark
